\begin{textsample}{POS Dim 3 – global\_south\_2025\_01 – Score 10.00 – global\_south\_2025\_01/120307907.txt}  \label{ex:f3_pos_404}
Harvard \textbf{University} has reached a significant settlement in \textbf{federal} court, agreeing to adopt the International Holocaust Remembrance Alliance ( IHRA ) definition of antisemitism.
This agreement is part of two legal cases filed by Jewish \textbf{students} who accused the \textbf{institution} of failing to protect them under Title VI of the 1964 Civil Rights Act, which prohibits discrimination based on race, color, or national origin at federally funded \textbf{institutions}.
The complaints arose from incidents of \textbf{harassment} and discrimination that \textbf{students} faced in the aftermath of the October 2023 Hamas attack on Israel and the subsequent Israeli military actions in Gaza. \textbf{Protests} on \textbf{campus}, which followed these events, fueled the grievances, with \textbf{students} claiming they were subjected to verbal and physical abuse, including bullying and intimidation.
They alleged that Harvard officials did not intervene effectively to stop the \textbf{harassment} from peers and \textbf{faculty} members.
The legal settlement represents a significant step for Harvard, as it aligns with a controversial definition of antisemitism established by the IHRA.
This definition describes antisemitism as"a certain @ @ @ @ @ @ @ @ @ @ toward Jews."It has drawn criticism, with opponents arguing that it could be used to suppress legitimate criticism of Israel.
The definition includes examples that could be seen as controversial, such as"claiming that the existence of a State of Israel is a racist endeavor"or applying"double standards"to Israel compared to other democratic nations.
Despite the opposition, the IHRA definition has been accepted by the US State Department and several European governments.
Pro-Israel advocacy groups have lobbied heavily for its adoption, seeing it as a necessary tool to combat antisemitism in modern times.
However, civil rights organizations have raised concerns, with the ACLU calling the definition"overbroad"and urging lawmakers to reject efforts to codify it.
The settlement is likely to attract significant attention from the Trump administration, which has been vocal about its stance on antisemitism and its interest in enforcing stricter policies on \textbf{campuses}.
In fact, this ruling comes after months of intense debates in the US over \textbf{campus} \textbf{protests} related to @ @ @ @ @ @ @ @ @ @, accusing \textbf{universities} of failing to protect Jewish \textbf{students} from \textbf{harassment}.
The controversy surrounding Harvard and other Ivy League \textbf{universities} has led to high-profile resignations.
Former Harvard President Claudine Gay stepped down earlier after facing criticism for not adequately addressing the \textbf{campus} \textbf{protests} and the alleged failure to protect Jewish \textbf{students}.
Similarly, leaders at Columbia \textbf{University} and the \textbf{University} of Pennsylvania faced similar pressures, with many conservatives accusing them of not doing enough to manage \textbf{protests}.
As part of the settlement, Harvard has agreed to publish an annual public report on violations of Title VI for the next five years, offering transparency on the \textbf{university} ' s handling of discrimination complaints.
The legal cases were initially dismissed by Harvard, but they were revived in 2024, with the court ruling that the cases should proceed.
The move by Harvard to adopt the IHRA definition of antisemitism and publish detailed reports is viewed as an attempt to address the concerns of Jewish \textbf{students} and comply with pressure from political leaders, including those in the Republican Party. @ @ @ @ @ @ @ @ @ @ broader discussions on free \textbf{speech}, \textbf{academic} freedom, and the definition of hate \textbf{speech} on \textbf{university} \textbf{campuses} in the future.

% matched lemmas: academic, campus, faculty, federal, harassment, institution, protest, speech, student, university
\end{textsample}
